\subsection{Small binary corners in a tree}

\begin{proposition}[Proved]
\label{CP2:A-tree-small-binary-corners}
Let $C$ be a finite connected category with at least two
objects and acyclic support, whose undirected comparability
graph is a tree. Suppose there is an object $v$ with
$|\operatorname{End}_C(v)|\leq11$, and every other corner
has order at most four. At each object $w\ne v$ whose
corner has order three or four, assume every nonempty
strict hom-set incident with $w$ has at most two arrows.
If $BC$ has bounded rational homology and
$\dim H^2(BC;\mathbb Q)\leq1$, every rational degree-two
product on $BC$ vanishes.
\end{proposition}
\begin{proof}
A tree support has no strict chain on three objects,
since its composite would supply the third edge of a
triangle. Thus each object is a source or sink.
Suppose, towards a contradiction, that a degree-two
product is nonzero. The second cohomology is then
one-dimensional, with a generator $u$ satisfying
$u^2\ne0$.

We describe a replacement at a small binary corner $w$,
preserving bounded rational homology, second Betti number
one, and a nonzero degree-two square. By opposition
suppose $w$ is a sink, and write $N=\operatorname{End}_C(w)$.
If the trivial right $\mathbb QN$-module is projective,
the actual orbit functor of
Lemma~\ref{CP2:A-high-orbits} replaces $w$ by a trivial
corner and all its incoming stalks by their $N$-orbit
sets. It is a rational homology equivalence and induces
a cohomology-ring isomorphism. Empty stalks remain empty
and nonempty stalks remain nonempty; cardinalities cannot
increase. All the required properties are preserved.

Otherwise the projectivity test and the order-four
minimal-ideal budget give a right-zero minimal band
$K$ with two or three elements, one row and trivial
maximal group. Each $k\in K$ is a right zero of the
whole monoid, so $BN$ is contractible. These are exactly
the alternatives used in
Lemma~\ref{CP2:A-high-small-binary}.
On each incoming stalk, all actors from $K$ are the
same idempotent map $e$. Indeed their images coincide;
on a set of cardinality at most two, either that common
image is a singleton, or all maps are identities.
On an empty stalk use its unique self-map. For $n\in N$,
the identities $nk=k$ and $kn\in K$ show
$ne=e=en$. Precomposition by any arrow of the opposite
block commutes with these actors, so this family is an
idempotent natural transformation of the entire connecting
profunctor, including its empty components.

Let $Y$ be its image subcategory. There are actual functors
\[
 R:C\longrightarrow Y,\qquad i:Y\hookrightarrow C,
 \qquad Ri=\operatorname{id}_Y.
\]
The relative mapping-cylinder argument in
Theorem~\ref{CP2:A-cut-retract} identifies the relative
quotient with a suspension, so positive relative
cohomology products are zero. If $H^2(BY;\mathbb Q)=0$,
then $u$ lifts to relative cohomology and $u^2=0$,
a contradiction. Otherwise the split injection
$R^*:H^2(BY;\mathbb Q)\to H^2(BC;\mathbb Q)$ is an
isomorphism, since its target has dimension one.
Choose $u_Y$ with $R^*u_Y=u$. Naturality gives
$R^*(u_Y^2)=u^2\ne0$, so the square remains nonzero.
The homology of $BY$ is a direct summand of that of
$BC$, hence remains bounded.

All $N$ actors are trivial on the image stalks of $Y$.
There is therefore an actual quotient functor replacing
$N$ by the trivial corner while keeping those stalks.
It is a homotopy equivalence: in the two-block model,
its corner map is $BN\to *$ and its bridge map is the
projection $BN\times T\to T$, while the opposite block
is fixed. Both displayed maps are homotopy equivalences
because $BN$ is contractible; the homotopy-pushout
comparison gives the assertion. This includes empty
stalks, since they contribute empty bar components.
The resulting category again has bounded homology,
second Betti number one, and a nonzero square. Its
objects and support are unchanged, and no stalk has
grown. The source case is the opposite construction.

Repeat this replacement at every $w\ne v$ of order
three or four. All future binary hypotheses persist,
because each step can only decrease incident cardinalities.
The final category has the same tree support, the corner
at $v$ still has order at most eleven, and every other
corner has order at most two. Proposition~\ref{CP2:A-tree-one-large-corner-eleven}
annihilates every degree-two product under boundedness,
contrary to the square preserved above. This proves
the result. Rank at most one also makes every mixed
product a scalar multiple of a square.
\end{proof}

\begin{lemma}[Proved: a trivial pendant with parallel arrows]
\label{CP2:A-trivial-pendant-parallel-circle}
Let $p$ be an object of a finite category with acyclic
support, strictly adjacent only to $q$. Suppose both
corner monoids are trivial and the nonempty hom-set
between them has cardinality $s\geq2$. Then
\[
 BC\simeq B(C\setminus\{p\})\vee
                 \bigvee^{s-1}S^1,
\]
where the wedge is taken at the component containing
$q$. In particular $H_1(BC;\mathbb Q)\ne0$.
\end{lemma}
\begin{proof}
The pendant condition makes the nerve the union of
the retained nerve and the full two-object nerve on
$p,q$, with intersection exactly the vertex $q$.
The two-object nerve is the connected graph with two
vertices and $s$ parallel edges. It contracts a spanning
edge to a wedge of $s-1$ circles, by a homotopy equivalence
relative to $q$. Attach this equivalence to the identity
of the retained nerve. This gives the stated wedge and
its direct summand $\mathbb Q^{s-1}$ in first homology.
\end{proof}
